% This file is intended to be inserted immediately after Appendix A.
% It contains exactly one section and uses the notation already introduced
% in the main text and in the cavity appendix.

\section{Central limit theorem for the overlap}
\label{sec:overlap-clt}

In this section we prove the Gaussian fluctuation theorem for the overlap.
We work at the endpoint $s=1$ of the smart path and suppress the subscript
$s$ from the notation.  Recall
\[
 Q_{ab}=R_{ab}-q,
 \qquad
 q=\mathbb E\tanh^2(h+\beta\sqrt q\,Z),
 \qquad
 r=\mathbb E\tanh^4(h+\beta\sqrt q\,Z),
\]
and put
\[
 a\coloneqq 1-2q+r,
 \qquad
 \alpha\coloneqq \beta^2a,
 \qquad
 \kappa\coloneqq1-4q+3r,
 \qquad
 \zeta\coloneqq2q+q^2-3r.
\]
Thus $\alpha$ is the Almeida--Thouless parameter.  Throughout this section
we assume
\[
 \alpha<1.
\]
By the argument preceding \eqref{eq:denominator-bounds},
\[
 \beta^2\kappa\le\alpha<1,
\]
and therefore
\begin{equation}
 1-\alpha>0,
 \qquad
 1-\beta^2\kappa>0.
 \label{eq:clt-positive-denominators}
\end{equation}

Define
\begin{equation}
 \sigma^2
 \coloneqq
 \frac{3a}{1-\alpha}
 -\frac{2\kappa}{1-\beta^2\kappa}
 -\frac{\zeta}{(1-\beta^2\kappa)^2}.
 \label{eq:clt-variance}
\end{equation}

\begin{theorem}[Central limit theorem for the overlap]
\label{thm:overlap-clt}
For every fixed $(\beta,h)$ in the strict Almeida--Thouless region,
\[
 \sqrt N\,(R_{12}-q)
 \xrightarrow{\mathrm d}
 \mathcal N(0,\sigma^2)
\]
with respect to the disorder-averaged two-replica Gibbs measure, where
$\sigma^2$ is given by \eqref{eq:clt-variance}.  Equivalently, for every
$t\in\mathbb R$,
\begin{align}
 \nu\!\left[
 \cos\!\left(t\sqrt N\,Q_{12}\right)
 \right]
 &\longrightarrow
 \exp\!\left(-\frac{\sigma^2t^2}{2}\right),
 \label{eq:clt-cos-limit}\\
 \nu\!\left[
 \sin\!\left(t\sqrt N\,Q_{12}\right)
 \right]
 &\longrightarrow 0.
 \label{eq:clt-sin-limit}
\end{align}
\end{theorem}

\begin{proof}
We divide the computation into several paragraphs.

\paragraph{Uniform second moment along the last-spin interpolation.}
Let $\nu_u$, $0\le u\le1$, be the last-spin interpolation from
Appendix~\ref{app:cavity-identities}, so that $\nu_1=\nu$.  Recall
\[
 \varepsilon_a=\sigma_N^a,
 \qquad
 Q^-_{ab}=\frac1N\sum_{i<N}\sigma_i^a\sigma_i^b-q,
 \qquad
 Q_{ab}=Q^-_{ab}+\frac1N\varepsilon_a\varepsilon_b.
\]
The optimal overlap estimate already proved in the main text gives
\begin{equation}
 \nu[Q_{12}^2]\le\frac{C}{N}.
 \label{eq:clt-endpoint-second-moment}
\end{equation}
Set
\[
 M_2(u)\coloneqq\nu_u[(Q^-_{12})^2].
\]
Applying \eqref{eq:cavity-derivative-proof} to
$F=(Q^-_{12})^2$, every term in $M_2'(u)$ has the form
\[
 \nu_u\!\left[
 (Q^-_{12})^2
 \varepsilon_a\varepsilon_bQ^-_{ab}
 \right].
\]
Since $|Q^-_{ab}|\le2$,
\[
 \left|
 \nu_u\!\left[
 (Q^-_{12})^2
 \varepsilon_a\varepsilon_bQ^-_{ab}
 \right]
 \right|
 \le2M_2(u).
\]
The number of terms in the cavity derivative is fixed, and $\beta$ is fixed,
hence
\begin{equation}
 |M_2'(u)|\le C M_2(u).
 \label{eq:clt-M2-differential}
\end{equation}
Moreover,
\[
 |Q^-_{12}|^2
 \le
 2|Q_{12}|^2+\frac2{N^2},
\]
so \eqref{eq:clt-endpoint-second-moment} gives
\[
 M_2(1)\le\frac{C}{N}.
\]
Gronwall's inequality applied backward from $u=1$ therefore yields
\begin{equation}
 \sup_{0\le u\le1}
 \nu_u[(Q^-_{12})^2]
 \le\frac{C}{N}.
 \label{eq:clt-uniform-cavity-second}
\end{equation}
By replica symmetry the same estimate holds for every fixed replica edge.
Consequently, for all fixed edges $ab,cd$,
\begin{align}
 \sup_{0\le u\le1}\nu_u[|Q^-_{ab}|]
 &\le\frac{C}{\sqrt N},
 \label{eq:clt-cavity-first}\\
 \sup_{0\le u\le1}
 \nu_u[|Q^-_{ab}Q^-_{cd}|]
 &\le\frac{C}{N}.
 \label{eq:clt-cavity-product}
\end{align}
The second inequality follows from Cauchy--Schwarz.

\paragraph{Three test-function correlations.}
Let $f\in C^2(\mathbb R)$ satisfy
\[
 \|f\|_\infty+\|f'\|_\infty+\|f''\|_\infty<\infty.
\]
Put
\[
 X_{ab}\coloneqq\sqrt N\,Q_{ab}.
\]
We consider the three replica-edge classes
\begin{align}
 T_A(f)&\coloneqq\nu[X_{12}f(X_{12})],
 \label{eq:clt-TA}\\
 T_B(f)&\coloneqq\nu[X_{13}f(X_{12})],
 \label{eq:clt-TB}\\
 T_C(f)&\coloneqq\nu[X_{34}f(X_{12})].
 \label{eq:clt-TC}
\end{align}
After the disorder average the sites are exchangeable.  Hence, for every
edge $ab$ appearing above,
\[
 \nu[Q_{ab}f(X_{12})]
 =
 \nu[(\varepsilon_a\varepsilon_b-q)f(X_{12})].
\]
Thus, if
\[
 \overline\varepsilon_{ab}
 \coloneqq
 \varepsilon_a\varepsilon_b-q,
\]
then
\begin{equation}
 T_E(f)
 =
 \sqrt N\,\nu[
 \overline\varepsilon_{e(E)}f(X_{12})],
 \qquad
 e(A)=12,\quad e(B)=13,\quad e(C)=34.
 \label{eq:clt-site-exchangeability}
\end{equation}

For $E\in\{A,B,C\}$ define
\[
 g_E(u)
 \coloneqq
 \sqrt N\,\nu_u[
 \overline\varepsilon_{e(E)}f(X_{12})].
\]
Then
\[
 g_E(1)=T_E(f).
\]

It is convenient to remove the last-spin contribution from the argument of
$f$.  Put
\begin{equation}
 Y
 \coloneqq
 \sqrt N\,Q^-_{12}+\frac q{\sqrt N}.
 \label{eq:clt-Yminus}
\end{equation}
Since
\[
 X_{12}
 =
 \sqrt N\,Q^-_{12}+\frac{\varepsilon_1\varepsilon_2}{\sqrt N},
\]
we have the exact identity
\begin{equation}
 X_{12}
 =
 Y+\frac{\overline\varepsilon_{12}}{\sqrt N}.
 \label{eq:clt-X-Y}
\end{equation}

At $u=0$ the last spin is independent of the first $N-1$ spins, and its
effective field has the distribution
\[
 h+\beta\sqrt q\,Z.
\]
Therefore, for distinct replica labels,
\[
 \mathbb E_0[\varepsilon_a\varepsilon_b]=q,
 \qquad
 \mathbb E_0[
 \varepsilon_a\varepsilon_b\varepsilon_c\varepsilon_d]=r.
\]
In particular,
\begin{align}
 \mathbb E_0[\overline\varepsilon_{12}^2]
 &=1-q^2,
 \label{eq:clt-source-A}\\
 \mathbb E_0[
 \overline\varepsilon_{13}\overline\varepsilon_{12}]
 &=q-q^2,
 \label{eq:clt-source-B}\\
 \mathbb E_0[
 \overline\varepsilon_{34}\overline\varepsilon_{12}]
 &=r-q^2.
 \label{eq:clt-source-C}
\end{align}
For example,
\[
 \mathbb E_0[
 \overline\varepsilon_{13}\overline\varepsilon_{12}]
 =
 \mathbb E_0[\varepsilon_2\varepsilon_3]
 -q\mathbb E_0[\varepsilon_1\varepsilon_3]
 -q\mathbb E_0[\varepsilon_1\varepsilon_2]
 +q^2
 =
 q-q^2.
\]

By Taylor's theorem and \eqref{eq:clt-X-Y},
\[
 f(X_{12})
 =
 f(Y)
 +\frac{\overline\varepsilon_{12}}{\sqrt N}f'(Y)
 +\rho_N,
\]
where
\[
 |\rho_N|
 \le
 \frac{\|f''\|_\infty}{2N}
 |\overline\varepsilon_{12}|^2
 \le
 \frac{2\|f''\|_\infty}{N}.
\]
Because
$\mathbb E_0[\overline\varepsilon_{e(E)}]=0$, independence at $u=0$
and \eqref{eq:clt-source-A}--\eqref{eq:clt-source-C} give
\begin{equation}
 g_E(0)
 =
 \theta_E\,\nu_0[f'(Y)]
 +O_f(N^{-1/2}),
 \label{eq:clt-g-zero-pre}
\end{equation}
where
\begin{equation}
 \theta_A=1-q^2,
 \qquad
 \theta_B=q-q^2,
 \qquad
 \theta_C=r-q^2.
 \label{eq:clt-theta}
\end{equation}
Here and below $O_f(\cdot)$ means that the implicit constant depends only on
$\beta,h$ and
$\|f\|_\infty+\|f'\|_\infty+\|f''\|_\infty$.

We next replace $\nu_0[f'(Y)]$ by the endpoint quantity
$\nu[f'(X_{12})]$.  Applying \eqref{eq:cavity-derivative-proof} to the
bounded cavity observable $f'(Y)$ and using
\eqref{eq:clt-cavity-first},
\[
 \left|
 \frac{\mathrm d}{\mathrm du}\nu_u[f'(Y)]
 \right|
 \le
 \frac{C\|f'\|_\infty}{\sqrt N}.
\]
Also, by \eqref{eq:clt-X-Y},
\[
 |f'(X_{12})-f'(Y)|
 \le
 \frac{2\|f''\|_\infty}{\sqrt N}.
\]
Hence
\begin{equation}
 \nu_0[f'(Y)]
 =
 \nu[f'(X_{12})]+O_f(N^{-1/2}).
 \label{eq:clt-derivative-endpoint-replacement}
\end{equation}
Combining \eqref{eq:clt-g-zero-pre} and
\eqref{eq:clt-derivative-endpoint-replacement},
\begin{equation}
 g_E(0)
 =
 \theta_E\,m_f+O_f(N^{-1/2}),
 \qquad
 m_f\coloneqq\nu[f'(X_{12})].
 \label{eq:clt-g-zero}
\end{equation}

\paragraph{First cavity derivative and the replica-edge matrix.}
Define the cavity analogues
\begin{align}
 \widetilde T_A(f)
 &\coloneqq
 \sqrt N\,\nu_0[Q^-_{12}f(Y)],
 \label{eq:clt-Ttilde-A}\\
 \widetilde T_B(f)
 &\coloneqq
 \sqrt N\,\nu_0[Q^-_{13}f(Y)],
 \label{eq:clt-Ttilde-B}\\
 \widetilde T_C(f)
 &\coloneqq
 \sqrt N\,\nu_0[Q^-_{34}f(Y)].
 \label{eq:clt-Ttilde-C}
\end{align}
We compute $g_E'(0)$ using
\eqref{eq:cavity-derivative-proof}.  In the derivative, replacing
$f(X_{12})$ by $f(Y)$ costs $O_f(N^{-1/2})$: indeed,
\[
 |f(X_{12})-f(Y)|
 \le
 \frac{2\|f'\|_\infty}{\sqrt N},
\]
and each derivative term contains one cavity overlap, whose first moment is
$O(N^{-1/2})$ by \eqref{eq:clt-cavity-first}.

Introduce
\[
 b_2\coloneqq\beta^2(1-q^2),
 \qquad
 b_1\coloneqq\beta^2(q-q^2),
 \qquad
 b_0\coloneqq\beta^2(r-q^2).
\]
The three rows of the cavity matrix are obtained as follows.

For $E=A$, use two replicas in
\eqref{eq:cavity-derivative-proof}.  The positive pair is $12$, the two
negative pairs are $13,23$, and the final pair is $34$.  By the edge rule,
\[
 g_A'(0)
 =
 b_2\widetilde T_A
 -4b_1\widetilde T_B
 +3b_0\widetilde T_C
 +O_f(N^{-1/2}).
 \label{eq:clt-row-A}
\]

For $E=B$, use three replicas.  The positive pairs $12,13,23$ contribute
\[
 b_1\widetilde T_A+(b_2+b_1)\widetilde T_B.
\]
The three negative terms are multiplied by $-3$.  The edges $14$ and $24$
give respectively the coefficients $b_1$ and $b_0$ in the shared-index
class, while $34$ gives the coefficient $b_1$ in the disjoint class.
Finally the coefficient of the new edge $45$ is $6b_0$.  Therefore
\begin{align}
 g_B'(0)
 &=
 b_1\widetilde T_A
 +(b_2+b_1)\widetilde T_B
 -3(b_1+b_0)\widetilde T_B
 -3b_1\widetilde T_C
 +6b_0\widetilde T_C
 +O_f(N^{-1/2})
 \nonumber\\
 &=
 b_1\widetilde T_A
 +(b_2-2b_1-3b_0)\widetilde T_B
 +(6b_0-3b_1)\widetilde T_C
 +O_f(N^{-1/2}).
 \label{eq:clt-row-B}
\end{align}

For $E=C$, use four replicas.  Among the six positive pairs, the edge $12$
is disjoint from $34$, the four edges $13,14,23,24$ share one endpoint with
$34$, and the edge $34$ is the same edge.  Hence the positive contribution is
\[
 b_0\widetilde T_A
 +4b_1\widetilde T_B
 +b_2\widetilde T_C.
\]
For the four negative terms, multiplied by $-4$, the edges $15,25$ are
disjoint from $34$ and belong to the shared-index class relative to $12$,
whereas $35,45$ share one endpoint with $34$ and belong to the disjoint
class relative to $12$.  Thus the negative contribution is
\[
 -8b_0\widetilde T_B-8b_1\widetilde T_C.
\]
The final edge $56$ is disjoint from $34$ and contributes
$10b_0\widetilde T_C$.  Therefore
\begin{equation}
 g_C'(0)
 =
 b_0\widetilde T_A
 +(4b_1-8b_0)\widetilde T_B
 +(b_2-8b_1+10b_0)\widetilde T_C
 +O_f(N^{-1/2}).
 \label{eq:clt-row-C}
\end{equation}

Thus, with
\[
 T(f)=
 \begin{pmatrix}
 T_A(f)\\ T_B(f)\\ T_C(f)
 \end{pmatrix},
 \qquad
 \widetilde T(f)=
 \begin{pmatrix}
 \widetilde T_A(f)\\
 \widetilde T_B(f)\\
 \widetilde T_C(f)
 \end{pmatrix},
 \qquad
 \theta=
 \begin{pmatrix}
 1-q^2\\ q-q^2\\ r-q^2
 \end{pmatrix},
\]
and
\begin{equation}
 M=
 \begin{pmatrix}
 b_2 & -4b_1 & 3b_0\\
 b_1 & b_2-2b_1-3b_0 & 6b_0-3b_1\\
 b_0 & 4b_1-8b_0 & b_2-8b_1+10b_0
 \end{pmatrix},
 \label{eq:clt-cavity-matrix}
\end{equation}
equations \eqref{eq:clt-row-A}--\eqref{eq:clt-row-C} become
\begin{equation}
 g'(0)
 =
 M\widetilde T(f)+O_f(N^{-1/2}).
 \label{eq:clt-g-prime-zero}
\end{equation}

\paragraph{Cavity Taylor remainder and endpoint replacement.}
Apply the cavity derivative formula once more.  Every term in $g_E''(u)$ is
a fixed numerical multiple of an expectation containing two cavity overlaps.
Since
\[
 \left|
 \sqrt N\,
 \overline\varepsilon_{e(E)}f(X_{12})
 \right|
 \le
 2\sqrt N\,\|f\|_\infty,
\]
equation \eqref{eq:clt-cavity-product} gives
\begin{equation}
 \sup_{0\le u\le1}|g_E''(u)|
 \le
 \frac{C_f}{\sqrt N}.
 \label{eq:clt-second-cavity-derivative}
\end{equation}
Taylor's formula in $u$ therefore yields
\begin{equation}
 T_E(f)=g_E(1)
 =
 g_E(0)+g_E'(0)+O_f(N^{-1/2}).
 \label{eq:clt-cavity-Taylor}
\end{equation}

It remains to replace $\widetilde T_E(f)$ by $T_E(f)$.  Consider, for
example,
\[
 H_A(u)
 =
 \sqrt N\,\nu_u[Q^-_{12}f(Y)].
\]
Differentiating with respect to $u$ produces one additional cavity overlap.
Using \eqref{eq:clt-cavity-product},
\[
 |H_A'(u)|
 \le
 \frac{C_f}{\sqrt N}.
\]
The same argument applies to the $B$ and $C$ classes.  Hence
\begin{equation}
 \sqrt N\,\nu_0[Q^-_{e(E)}f(Y)]
 =
 \sqrt N\,\nu[Q^-_{e(E)}f(Y)]
 +O_f(N^{-1/2}).
 \label{eq:clt-u-endpoint-replacement}
\end{equation}
At $u=1$,
\[
 Q_{e(E)}-Q^-_{e(E)}
 =
 \frac1N\varepsilon_{e(E)},
\]
so
\[
 \sqrt N\,
 \left|
 \nu[(Q_{e(E)}-Q^-_{e(E)})f(X_{12})]
 \right|
 \le
 \frac{\|f\|_\infty}{\sqrt N}.
\]
Furthermore,
\[
 |f(X_{12})-f(Y)|
 \le
 \frac{2\|f'\|_\infty}{\sqrt N},
\]
and therefore, by \eqref{eq:clt-cavity-first},
\[
 \sqrt N\,
 \nu\!\left[
 |Q^-_{e(E)}|
 |f(X_{12})-f(Y)|
 \right]
 \le
 \frac{C_f}{\sqrt N}.
\]
Combining these estimates,
\begin{equation}
 \widetilde T_E(f)
 =
 T_E(f)+O_f(N^{-1/2}),
 \qquad E\in\{A,B,C\}.
 \label{eq:clt-Ttilde-T}
\end{equation}

Equations \eqref{eq:clt-g-zero},
\eqref{eq:clt-g-prime-zero},
\eqref{eq:clt-cavity-Taylor}, and
\eqref{eq:clt-Ttilde-T} give the approximate Stein system
\begin{equation}
 T(f)
 =
 \theta\,m_f
 +
 M T(f)
 +
 O_f(N^{-1/2}),
 \qquad
 m_f=\nu[f'(X_{12})].
 \label{eq:clt-vector-Stein}
\end{equation}

\paragraph{Diagonalization in the three cavity modes.}
For the test-function correlations define
\begin{align}
 U_f&\coloneqq T_A(f)-4T_B(f)+3T_C(f),
 \label{eq:clt-Uf}\\
 V_f&\coloneqq2T_B(f)-3T_C(f),
 \label{eq:clt-Vf}\\
 D_f&\coloneqq T_A(f)-2T_B(f)+T_C(f).
 \label{eq:clt-Df}
\end{align}
The same deterministic change of basis as in the cavity closure triangularizes
the matrix \eqref{eq:clt-cavity-matrix}.  We record the coefficient
calculation explicitly.

First,
\begin{align*}
 \theta_A-4\theta_B+3\theta_C
 &=
 (1-q^2)-4(q-q^2)+3(r-q^2)\\
 &=1-4q+3r
 =\kappa,
\\
 2\theta_B-3\theta_C
 &=
 2(q-q^2)-3(r-q^2)\\
 &=2q+q^2-3r
 =\zeta,
\\
 \theta_A-2\theta_B+\theta_C
 &=
 (1-q^2)-2(q-q^2)+(r-q^2)\\
 &=1-2q+r
 =a.
\end{align*}
Second, direct multiplication of the rows in
\eqref{eq:clt-Uf}--\eqref{eq:clt-Df} by $M$ gives
\begin{align*}
 U(MT)
 &=\beta^2\kappa\,U_f,\\
 V(MT)
 &=\beta^2\zeta\,U_f+\beta^2\kappa\,V_f,\\
 D(MT)
 &=\beta^2a\,D_f
 =\alpha D_f.
\end{align*}
Thus \eqref{eq:clt-vector-Stein} is equivalent to
\begin{align}
 (1-\beta^2\kappa)U_f
 &=
 \kappa\,m_f+O_f(N^{-1/2}),
 \label{eq:clt-Stein-U}\\
 (1-\beta^2\kappa)V_f-\beta^2\zeta U_f
 &=
 \zeta\,m_f+O_f(N^{-1/2}),
 \label{eq:clt-Stein-V}\\
 (1-\alpha)D_f
 &=
 a\,m_f+O_f(N^{-1/2}).
 \label{eq:clt-Stein-D}
\end{align}
Using \eqref{eq:clt-positive-denominators},
\begin{align}
 U_f
 &=
 \frac{\kappa}{1-\beta^2\kappa}\,m_f
 +O_f(N^{-1/2}),
 \label{eq:clt-U-solved}\\
 D_f
 &=
 \frac{a}{1-\alpha}\,m_f
 +O_f(N^{-1/2}).
 \label{eq:clt-D-solved}
\end{align}
Substituting \eqref{eq:clt-U-solved} into
\eqref{eq:clt-Stein-V},
\begin{align*}
 V_f
 &=
 \frac{
 \zeta m_f+\beta^2\zeta U_f
 }{1-\beta^2\kappa}
 +O_f(N^{-1/2})
\\
 &=
 \frac{\zeta}{1-\beta^2\kappa}
 \left(
 1+
 \frac{\beta^2\kappa}{1-\beta^2\kappa}
 \right)m_f
 +O_f(N^{-1/2})
\\
 &=
 \frac{\zeta}{(1-\beta^2\kappa)^2}\,m_f
 +O_f(N^{-1/2}).
 \label{eq:clt-V-solved}
\end{align*}
The inverse mode transformation is
\begin{equation}
 T_A(f)=-V_f-2U_f+3D_f.
 \label{eq:clt-inverse-mode-A}
\end{equation}
Therefore
\begin{align}
 \nu[X_{12}f(X_{12})]
 &=
 \left(
 \frac{3a}{1-\alpha}
 -\frac{2\kappa}{1-\beta^2\kappa}
 -\frac{\zeta}{(1-\beta^2\kappa)^2}
 \right)
 \nu[f'(X_{12})]
 +O_f(N^{-1/2})
 \nonumber\\
 &=
 \sigma^2\,\nu[f'(X_{12})]
 +O_f(N^{-1/2}).
 \label{eq:clt-scalar-Stein}
\end{align}
This is the required approximate Gaussian integration-by-parts identity.

Notice that the choice $f\equiv1$ gives
\begin{equation}
 \nu[X_{12}]
 =
 O(N^{-1/2}),
 \label{eq:clt-bias}
\end{equation}
so the centering by the infinite-volume parameter $q$ has no residual
order-one bias.

\paragraph{Identification and nonnegativity of the variance.}
For completeness, we check that \eqref{eq:clt-variance} is the limiting
second moment.  From the cavity equations at $s=1$, the already proved
estimate
\[
 \nu[|Q_{12}|^3]=o(N^{-1})
\]
implies that the cavity remainders are $o(N^{-1})$.  Therefore
\begin{equation}
 NU
 \longrightarrow
 \frac{\kappa}{1-\beta^2\kappa}.
 \label{eq:clt-second-U}
\end{equation}
Next,
\[
 (1-\beta^2\kappa)NV
 -
 \beta^2\zeta\,NU
 =
 \zeta+o(1),
\]
and \eqref{eq:clt-second-U} gives
\begin{align}
 NV
 &\longrightarrow
 \frac{
 \zeta+
 \beta^2\zeta\kappa/(1-\beta^2\kappa)
 }{1-\beta^2\kappa}
 \nonumber\\
 &=
 \frac{\zeta}{(1-\beta^2\kappa)^2}.
 \label{eq:clt-second-V}
\end{align}
Likewise,
\begin{equation}
 ND
 \longrightarrow
 \frac{a}{1-\alpha}.
 \label{eq:clt-second-D}
\end{equation}
Since
\[
 A=-V-2U+3D,
\]
we conclude that
\begin{align}
 N\nu[Q_{12}^2]
 =NA
 &\longrightarrow
 -\frac{\zeta}{(1-\beta^2\kappa)^2}
 -\frac{2\kappa}{1-\beta^2\kappa}
 +\frac{3a}{1-\alpha}
 \nonumber\\
 &=\sigma^2.
 \label{eq:clt-second-moment-limit}
\end{align}
In particular,
\begin{equation}
 \sigma^2\ge0.
 \label{eq:clt-variance-nonnegative}
\end{equation}

\paragraph{Characteristic-function ODE.}
Define
\begin{align}
 C_N(t)
 &\coloneqq
 \nu[\cos(tX_{12})],
 \label{eq:clt-CN}\\
 S_N(t)
 &\coloneqq
 \nu[\sin(tX_{12})].
 \label{eq:clt-SN}
\end{align}
For fixed $N$ these functions are differentiable and
\begin{align}
 C_N'(t)
 &=
 -\nu[X_{12}\sin(tX_{12})],
 \label{eq:clt-CN-derivative}\\
 S_N'(t)
 &=
 \nu[X_{12}\cos(tX_{12})].
 \label{eq:clt-SN-derivative}
\end{align}
Fix $T>0$.  For $|t|\le T$, the functions
\[
 f_t(x)=\sin(tx),
 \qquad
 \widetilde f_t(x)=\cos(tx)
\]
satisfy
\[
 \|f_t\|_\infty,\|\widetilde f_t\|_\infty\le1,
 \qquad
 \|f_t'\|_\infty,\|\widetilde f_t'\|_\infty\le T,
 \qquad
 \|f_t''\|_\infty,\|\widetilde f_t''\|_\infty\le T^2.
\]
Hence the error in \eqref{eq:clt-scalar-Stein} is uniform for
$|t|\le T$.

Applying \eqref{eq:clt-scalar-Stein} to $f_t$ gives
\[
 \nu[X_{12}\sin(tX_{12})]
 =
 \sigma^2t\,C_N(t)+O_T(N^{-1/2}),
\]
and therefore
\begin{equation}
 C_N'(t)+\sigma^2tC_N(t)
 =
 O_T(N^{-1/2})
 \qquad(|t|\le T).
 \label{eq:clt-C-approx-ODE}
\end{equation}
Applying \eqref{eq:clt-scalar-Stein} to $\widetilde f_t$ gives
\[
 \nu[X_{12}\cos(tX_{12})]
 =
 -\sigma^2t\,S_N(t)+O_T(N^{-1/2}),
\]
and hence
\begin{equation}
 S_N'(t)+\sigma^2tS_N(t)
 =
 O_T(N^{-1/2})
 \qquad(|t|\le T).
 \label{eq:clt-S-approx-ODE}
\end{equation}

Multiply \eqref{eq:clt-C-approx-ODE} by
$\exp(\sigma^2t^2/2)$.  Since
\[
 C_N(0)=1,
\]
we obtain, for $|t|\le T$,
\begin{align*}
 \left|
 e^{\sigma^2t^2/2}C_N(t)-1
 \right|
 &\le
 \int_0^{|t|}
 e^{\sigma^2u^2/2}
 \frac{C_T}{\sqrt N}\,\mathrm du\\
 &\le
 \frac{C_TT}{\sqrt N}
 e^{\sigma^2T^2/2}.
\end{align*}
Therefore
\begin{equation}
 \sup_{|t|\le T}
 \left|
 C_N(t)-e^{-\sigma^2t^2/2}
 \right|
 \longrightarrow0.
 \label{eq:clt-C-uniform}
\end{equation}
Similarly, since $S_N(0)=0$,
\begin{align*}
 \left|
 e^{\sigma^2t^2/2}S_N(t)
 \right|
 &\le
 \int_0^{|t|}
 e^{\sigma^2u^2/2}
 \frac{C_T}{\sqrt N}\,\mathrm du,
\end{align*}
so
\begin{equation}
 \sup_{|t|\le T}|S_N(t)|
 \longrightarrow0.
 \label{eq:clt-S-uniform}
\end{equation}
Since $T>0$ was arbitrary,
\eqref{eq:clt-C-uniform} and \eqref{eq:clt-S-uniform} prove
\eqref{eq:clt-cos-limit} and \eqref{eq:clt-sin-limit}.  Equivalently,
\[
 \nu[e^{itX_{12}}]
 \longrightarrow
 e^{-\sigma^2t^2/2}
 \qquad(t\in\mathbb R).
\]
The right-hand side is the characteristic function of
$\mathcal N(0,\sigma^2)$ by
\eqref{eq:clt-variance-nonnegative}.  The continuity theorem for
characteristic functions now gives
\[
 X_{12}
 =
 \sqrt N\,(R_{12}-q)
 \xrightarrow{\mathrm d}
 \mathcal N(0,\sigma^2),
\]
which completes the proof.
\end{proof}
